\documentclass[10pt, a4paper]{report}
% IEEE body text = 10pt. Using report class for chapter support.

% =============================================================================
% PACKAGES
% =============================================================================

% --- Encoding & Language ---
\usepackage[utf8]{inputenc}
\usepackage[T1]{fontenc}
\usepackage[british]{babel}

% --- Typography ---
% IEEE uses Times Roman throughout
\usepackage{mathptmx}              % Times Roman for text + math
\usepackage[scaled=0.90]{helvet}   % Helvetica (sans-serif fallback)
\usepackage{courier}               % Courier for monospace / code

% --- Page Geometry (IEEE A4 spec) ---
% IEEE A4: side = 13mm, top = 19mm, bottom = 43mm
% Single-column width = 184mm
\usepackage[
    a4paper,
    left=13mm,
    right=13mm,
    top=19mm,
    bottom=43mm,
    headheight=12pt,
    headsep=10pt,
    footskip=20pt
]{geometry}

% --- Line Spacing ---
% IEEE body text: 10pt font / 12pt line spacing = factor of 1.2
% LaTeX default for 10pt is ~12pt baselineskip, so single spacing is correct.
\usepackage{setspace}
\singlespacing

% --- Headers & Footers ---
\usepackage{fancyhdr}
\pagestyle{fancy}
\fancyhf{}
\fancyhead[L]{\footnotesize\nouppercase{\leftmark}}
\fancyhead[R]{\footnotesize\thepage}
\renewcommand{\headrulewidth}{0.4pt}
\fancypagestyle{plain}{%
    \fancyhf{}
    \fancyfoot[C]{\footnotesize\thepage}
    \renewcommand{\headrulewidth}{0pt}
}

% --- Graphics & Figures ---
\usepackage{graphicx}
\usepackage{float}
\usepackage{subcaption}

% IEEE Figure captions: 8pt, 9pt line spacing, centred below figure
% IEEE Table captions: 8pt, 9pt line spacing, centred above table, small caps label
\usepackage[
    font={footnotesize, singlespacing},  % 8pt at 10pt base
    labelfont=bf,
    labelsep=period,
    justification=centering,
    tableposition=top,
    figureposition=bottom
]{caption}

% Override table caption label to use small caps "TABLE" with Roman numerals
\DeclareCaptionLabelFormat{ieeetable}{\textsc{Table~#2}}
\captionsetup[table]{labelformat=ieeetable, labelsep=newline, justification=centering}
\renewcommand{\thetable}{\Roman{table}}

% Figure numbering: Arabic, label "Fig."
\DeclareCaptionLabelFormat{ieeefigure}{Fig.~#2}
\captionsetup[figure]{labelformat=ieeefigure, labelsep=period, justification=centering}

% --- Tables ---
\usepackage{booktabs}
\usepackage{array}
\usepackage{multirow}
\usepackage{tabularx}

% --- Mathematics ---
\usepackage{amsmath}
\usepackage{amssymb}
\usepackage{bm}
\usepackage{braket}

% --- Colours ---
\usepackage{xcolor}
\definecolor{codeBackground}{RGB}{248,248,250}
\definecolor{codeFrame}{RGB}{190,190,200}
\definecolor{codeKeyword}{RGB}{0,80,160}
\definecolor{codeComment}{RGB}{80,128,80}
\definecolor{codeString}{RGB}{160,48,48}
\definecolor{codeNumber}{RGB}{140,140,150}
\definecolor{linkBlue}{RGB}{0,60,120}

% --- Code Listings ---
\usepackage{listings}
\lstdefinestyle{ieeeCode}{
    backgroundcolor=\color{codeBackground},
    frame=single,
    rulecolor=\color{codeFrame},
    framesep=3pt,
    xleftmargin=6pt,
    xrightmargin=4pt,
    % IEEE code should be small --- 8pt matches footnote/caption size
    basicstyle=\ttfamily\scriptsize\linespread{1.0}\selectfont,
    keywordstyle=\color{codeKeyword}\bfseries,
    commentstyle=\color{codeComment}\itshape,
    stringstyle=\color{codeString},
    numberstyle=\tiny\color{codeNumber}\ttfamily,
    numbers=left,
    numbersep=6pt,
    breaklines=true,
    breakatwhitespace=false,
    showstringspaces=false,
    captionpos=b,
    tabsize=4,
    columns=flexible,
    keepspaces=true,
    aboveskip=10pt,
    belowskip=6pt,
    abovecaptionskip=4pt,
    belowcaptionskip=2pt,
    literate={-}{-}1
}
\lstset{style=ieeeCode}
\renewcommand{\lstlistingname}{Listing}
\renewcommand{\lstlistlistingname}{List of Listings}

% --- Algorithms ---
\usepackage[ruled, vlined, linesnumbered]{algorithm2e}
\SetAlgoLined
\SetKwComment{Comment}{$\triangleright$\ }{}
\SetAlCapSkip{6pt}
% Match algorithm caption size to IEEE caption spec (8pt)
\SetAlCapFnt{\footnotesize}
\SetAlCapNameFnt{\footnotesize\bfseries}

% --- Hyperlinks ---
\usepackage[
    colorlinks=true,
    linkcolor=black,
    citecolor=linkBlue,
    urlcolor=linkBlue,
    bookmarks=true,
    bookmarksnumbered=true,
    pdfauthor={Om Kalyan},
    pdftitle={Beyond Classical Cryptography: Simulating and Benchmarking QKD},
    pdfsubject={BEng Dissertation},
    pdfkeywords={QKD, BB84, B92, E91, Qiskit, Post-Quantum Cryptography}
]{hyperref}
\usepackage[noabbrev, capitalise]{cleveref}

% --- Bibliography ---
% IEEE references: 8pt, 9pt line spacing, bracketed Arabic numerals
\usepackage{cite}

% --- Lists ---
\usepackage{enumitem}
\setlist{nosep, leftmargin=*}

% --- Paragraph formatting ---
% IEEE uses indented paragraphs with no extra inter-paragraph space
\setlength{\parindent}{1em}
\setlength{\parskip}{0pt}

% =============================================================================
% HEADING STYLES --- IEEE Conference Manuscript Format Definition
% =============================================================================
\usepackage{titlesec}

% --- CHAPTER ─────────────────────────────────────────────────────────────────
% Dissertation addition (IEEE conf. papers have no chapters).
% Styled as a prominent centred heading, consistent with IEEE aesthetic.
% Uses 14pt bold small caps, centred, with "CHAPTER I" label above.
\titleformat{\chapter}[display]
    {\normalfont\fontsize{14}{16}\bfseries\scshape\centering}
    {\MakeUppercase{\chaptertitlename}~\thechapter}
    {6pt}
    {\fontsize{14}{16}\selectfont\bfseries}
\titlespacing*{\chapter}{0pt}{-10pt}{20pt}

% --- SECTION (IEEE "Section Heads (1)") ──────────────────────────────────────
% Spec: 10pt regular, 14pt line spacing, centred,
%       Roman numerals (I., II., ...), CAPITALS / SMALL CAPS
\renewcommand{\thesection}{\Roman{section}}
\titleformat{\section}
    {\normalfont\fontsize{10}{14}\selectfont\scshape\centering}
    {\thesection.}{0.5em}{}
\titlespacing*{\section}{0pt}{14pt}{8pt}

% --- SUBSECTION (IEEE "Heading 2") ──────────────────────────────────────────
% Spec: 10pt italic, 12pt line spacing, left-justified,
%       Capital alpha numerals (A., B., ...)
\renewcommand{\thesubsection}{\Alph{subsection}}
\titleformat{\subsection}
    {\normalfont\fontsize{10}{12}\selectfont\itshape}
    {\thesubsection.}{0.5em}{}
\titlespacing*{\subsection}{0pt}{10pt}{4pt}

% --- SUBSUBSECTION (IEEE "Heading 3") ───────────────────────────────────────
% Spec: 10pt italic, 12pt line spacing, left-justified,
%       Arabic numeral + right paren, run-in, indent 1em
\renewcommand{\thesubsubsection}{\arabic{subsubsection}}
\titleformat{\subsubsection}[runin]
    {\normalfont\fontsize{10}{12}\selectfont\itshape}
    {\thesubsubsection)}{0.5em}{}[:]
\titlespacing*{\subsubsection}{1em}{8pt}{0.5em}

% --- PARAGRAPH (IEEE "Heading 4") ───────────────────────────────────────────
% Spec: 10pt italic, left-justified,
%       Lowercase alpha + right paren + colon, run-in, indent 2em
\renewcommand{\theparagraph}{\alph{paragraph}}
\titleformat{\paragraph}[runin]
    {\normalfont\fontsize{10}{12}\selectfont\itshape}
    {\theparagraph)}{0.5em}{}[:]
\titlespacing*{\paragraph}{2em}{6pt}{0.5em}

% Enable numbering depth through paragraph level
\setcounter{secnumdepth}{4}
\setcounter{tocdepth}{3}

% =============================================================================
% CUSTOM COMMANDS (QKD Notation)
% =============================================================================

\newcommand{\QBER}{\ensuremath{\mathrm{QBER}}}
\newcommand{\keyrate}{\ensuremath{R_{\mathrm{sec}}}}

% =============================================================================
% FOOTNOTE SIZE (IEEE: 8pt, 9pt line spacing)
% =============================================================================
\renewcommand{\footnotesize}{\fontsize{8}{9}\selectfont}

% =============================================================================
% DOCUMENT
% =============================================================================
\begin{document}

% ─────────────────────────────────────────────────────────────────────────────
% TITLE PAGE
% ─────────────────────────────────────────────────────────────────────────────
\begin{titlepage}
    \centering
    \thispagestyle{empty}
    \vspace*{3cm}

    % --- Title: IEEE spec = 24pt regular, 28pt line spacing, centred ---
    {\fontsize{28}{32}\selectfont
    Beyond Classical Cryptography:\par
    Simulating \& Benchmarking\par
    Quantum Key Distribution\par}

    \vspace{1.0cm}

    % --- Author Name: IEEE spec = 11pt regular, 12pt line spacing, centred ---
    {\fontsize{14}{15}\selectfont
    Om Kalyan\par}

    \vspace{0.5cm}

    {\fontsize{10}{11}\selectfont
    230733440\par}

    \vspace{1.0cm}

    % --- Author Affiliation: IEEE spec = 10pt regular, 11pt line spacing, centred ---
    {\fontsize{10}{11}\selectfont
    A dissertation submitted in partial fulfilment of the requirements\\
    for the degree of\par}

    \vspace{0.4cm}

    {\fontsize{10}{11}\selectfont\bfseries
    Bachelor of Engineering\\
    Electronics \& Computer Engineering\par}

    \vspace{1.0cm}

    % --- University Crest ---
    
\includegraphics[width=0.24\textwidth]{misc images/ncl logo vector.png}

    \vspace{1.0cm}

    % --- Institution: 11pt to match author name tier ---
    {\fontsize{11}{13}\selectfont
    School of Engineering\\
    Newcastle University\par}

    \vspace{0.6cm}

    % --- Date ---
    {\fontsize{10}{11}\selectfont 2026\par}

\end{titlepage}

% ─────────────────────────────────────────────────────────────────────────────
% FRONT MATTER
% ─────────────────────────────────────────────────────────────────────────────
\pagenumbering{roman}

% --- Abstract ---
\clearpage
\chapter*{Abstract}
\addcontentsline{toc}{chapter}{Abstract}

% IEEE Abstract: 9pt bold, "Abstract---" bold italic, 10pt line spacing
{\fontsize{9}{10}\selectfont\bfseries
Quantum algorithms such as Shor's factorisation threaten public-key cryptography systems like RSA and ECC, which rely on computational hardness for security. Quantum Key Distribution (QKD) offers a fundamentally different security paradigm: rather than relying on computational complexity, QKD derives information-theoretic security from the laws of quantum mechanics, ensuring privacy regardless of adversarial computing power. This project presents a simulation and benchmarking platform implementing three QKD protocols: BB84, B92, and E91 in Python3 using the Qiskit SDK. The platform compares Prepare \& Measure approaches against Entanglement-Based key distribution, evaluating protocol performance first under ideal channel conditions to establish theoretical baselines, then under realistic constraints including eavesdropper presence, quantum bit error rates, and channel losses.\par
}

% --- Acknowledgements ---
\clearpage
\chapter*{Acknowledgements}
\addcontentsline{toc}{chapter}{Acknowledgements}

I would like to thank my supervisors Martin Johnston and Oktay \c{C}etinkaya for their guidance\ldots

% --- Table of Contents ---
\clearpage
\tableofcontents

% --- List of Figures ---
\clearpage
\listoffigures
\addcontentsline{toc}{chapter}{List of Figures}

% --- List of Tables ---
\clearpage
\listoftables
\addcontentsline{toc}{chapter}{List of Tables}

% --- List of Listings ---
\clearpage
\lstlistoflistings
\addcontentsline{toc}{chapter}{List of Listings}

% --- List of Algorithms ---
\clearpage
\listofalgorithms
\addcontentsline{toc}{chapter}{List of Algorithms}

% ─────────────────────────────────────────────────────────────────────────────
% MAIN CONTENT
% ─────────────────────────────────────────────────────────────────────────────
\clearpage
\pagenumbering{arabic}

% ═══════════════════════════════════════════════════════════════════════════════
\chapter{Introduction}
% ═══════════════════════════════════════════════════════════════════════════════

\section{A Brief History of Cryptology}

% TODO: A brief introduction to QKD --- why it is needed, who benefits
% from the technology, and where it is deployed in the real world.

% this should talk abput classical cryptography first. an overview of why we need secure systems
Since the dawn of days, humans have sought ways to communicate in secret. It only takes a glance back at the ancient societies of Mesopotamia, India, and China to understand where the foundations of cryptology were laid some 3000 years ago \cite{Kahn1967Code}. Whether it be the early use of Substitution or Transposition schemes, these early cryptographic systems established the core principle that would dominate cryptology for millennia: security through the obscurity of the method itself. 

Back in the 9th century, humanity experienced the first major analytical breakthrough when the Arab polymath Al-Kindi developed frequency analysis, a technique of cryptanalysis capable of breaking any monoalphabetic substitution cipher of the age by exploiting the statistical patterns of language itself \cite{AlKadi1992}. This initiated a centuries-long cycle of invention and defeat that defined the field of cryptology, with newly designed ciphers broken by increasingly sophisticated cryptanalysis. Fast-forward to the 16th century, Blaise de Vigen\`ere built on the work of Giovan Bellaso to introduce a polyalphabetic cipher that would go on to secure the world for 300 years, until the great minds of Babbage and Kasiski developed statistical methods to compromise the scheme in the 19th century \cite{Singh1999, Bauer2021}.

The dawn of the Second World War transformed cryptology from a niche discipline into a strategic imperative. Repeating this "Tug of War", Germany's Engima, a rotor-based electromechanical cipher, was famously broken at Bletchley Park by the Father of Theoretical Computer Science - Alan Turing, and his colleagues \cite{Hodges2014}. Despite this technological revolution, security through simplicity was demonstrated in the Pacific theatre by the US Marine Corps' deployment of Navajo Code Talkers who used their native language as the basis for a battlefield cipher. A truly innovative system that rested security on the extreme phonetic distance between the Navajo and Japanese languages remained unbroken until its declassification in 1968 \cite{Huffman2000}. 

Looking to the post-war era, Claude Shannon redefined cryptography using rigorous mathematical foundations in his initially classified 1945 paper on secrecy systems, introducing concepts such as confusion, diffusion, and perfect secrecy \cite{Shannon1945}. This revolutionary framework, later publicly published in 1948, fueled the development of the modern cryptographic infrastructure we rely on today: the Data Encryption Standard (DES) in 1977, the RSA public-key cryptosystem in 1978, and the Advanced Encryption Standard (AES) in 2001 \cite{Shannon1948, DiffieHellman1976, RSA1978}. Peeling back the layers of essentially all the digital communications we use today, from online banking and e-commerce to government and military networks, we find these algorithms underpinning the security on which we rely. 

\section{The Quantum Threat to Classical Cryptography}

% TODO: Shor's algorithm, the quantum computing roadmap, and the need for
% information-theoretically secure key exchange.
% new intro para
The security of modern public-key cryptography rests on the computational hardness of problems like integer factorisation and the discrete logarithm. Shor's algorithm \cite{Shor1997} demonstrated that a sufficiently powerful quantum computer could solve both in polynomial time, rendering schemes such as RSA and ECC fundamentally insecure. Two complementary responses have emerged: Post-Quantum Cryptography (PQC), which replaces vulnerable algorithms with quantum-resistant mathematical problems, and Quantum Key Distribution (QKD), which exploits the laws of quantum mechanics to distribute encryption keys with information-theoretic security. This project focuses on the latter, building a simulation and benchmarking platform for QKD protocols including BB84, B92, and E91.


\newpage
\section{Aims \& Objectives}

% TODO: Outline the project aims and objectives.
% Describe the overall research message and trajectory.
The aim of this project is to design, implement, and evaluate a modular simulation and benchmarking platform for Discrete-Variable Quantum Key Distribution protocols, enabling quantitative comparison of protocol performance under realistic channel conditions, eavesdropping attacks, and post-processing constraints. All simulations are conducted in Python using IBM's Qiskit framework and the Aer noise simulator. The specific objectives are organised into six categories.

\subsection{Literature Review \& Analytical Framework}

\begin{enumerate}[label=\textbf{O1.\arabic*}, leftmargin=1.5cm]
    \item Conduct a critical review of QKD protocol literature, covering the theoretical foundations of BB84 \cite{BB84}, B92 \cite{B92}, and E91 \cite{E91}, their respective security proofs, and known practical limitations. This review will establish the analytical framework and parameter choices (noise models, error correction efficiency, security bounds) used throughout the project.

    \item Identify and document the key metrics by which QKD protocols are evaluated in the literature, such as: Quantum Bit Error Rate (QBER), sifted key rate, secure key rate, maximum tolerable error threshold, and sifting efficiency. These metrics will form the basis of all subsequent benchmarking.
\end{enumerate}

\subsection{Protocol Implementation}

\begin{enumerate}[label=\textbf{O2.\arabic*}, leftmargin=1.5cm]
    \item \textbf{BB84 (Four-State, Two-Basis):} Implement the BB84 protocol as the primary baseline simulation. The implementation models Alice's random bit and basis selection, qubit state preparation in the rectilinear ($Z$) or diagonal ($X$) basis, quantum transmission through a configurable noisy channel, Bob's random basis measurement, and public basis reconciliation (sifting). The implementation is considered complete when the sifting efficiency converges to the theoretically expected $\sim$50\% across repeated simulation runs. % refine with simualtions

    \item \textbf{B92 (Two Non-Orthogonal States):} Implement Bennett's 1992 protocol as the primary comparative prepare and measure protocol, encoding bit values using only two non-orthogonal states ($|0\rangle$ and $|+\rangle$). Bob performs a measurement that can yield a conclusive result, an inconclusive result, or a photon loss event. The implementation is validated by confirming that the sifting efficiency converges to $\sim$25\% and that the QBER threshold falls below that of BB84, consistent with B92's reduced noise tolerance.

    \item \textbf{E91 (Entanglement-Based):} Implement Ekert's 1991 protocol. This models the generation and distribution of maximally entangled Bell pairs, independent random basis measurements from three orientations per party, key extraction from matching-basis outcomes, and eavesdropper detection via the CHSH inequality. The implementation is validated by confirming that the computed CHSH parameter $|S|$ converges to $2\sqrt{2} \approx 2.828$ in the absence of eavesdropping, and degrades towards the classical bound of 2 under increasing interception rate.
\end{enumerate}

\subsection{Channel Modelling and Eavesdropper Simulation}

\begin{enumerate}[label=\textbf{O3.\arabic*}, leftmargin=1.5cm]
    \item Model quantum channel effects using Qiskit Aer's noise framework: depolarising noise, and phase damping. Noise is applied exclusively to the quantum channel via an identity gate channel marker, not to local preparation or measurement operations, ensuring the simulation accurately represents a protocol where Alice's and Bob's local devices are trusted.
    
    \item Implement an intercept-resend eavesdropping attack with a configurable interception probability $p \in [0,1]$. Eve intercepts each qubit with probability $p$, measures in a randomly chosen basis, and re-transmits to Bob. Sweep the interception rate from 0\% to 100\%  record the resulting QBER for each protocol.
    
\end{enumerate}

\subsection{Post-Processing and Secure Key Rate Analysis}

\begin{enumerate}[label=\textbf{O4.\arabic*}, leftmargin=1.5cm]
    \item Calculate the secure key rate using the GLLP formula \cite{GLLP2004}, incorporating the error correction efficiency factor $f_{EC} = 1.16$ corresponding to the CASCADE protocol \cite{Cascade1994}. This yields a practical BB84 QBER security threshold of $\sim$9.8\%, compared to the idealised 11\% Shor-Preskill bound at $f_{EC} = 1.0$ \cite{ShorPreskill2000}. Document and justify this distinction, as it has direct engineering significance for determining whether a given fibre link can sustain secure key generation.

    \item Compute the mutual information quantities $I(A{:}B)$ (Alice-Bob) and $I(A{:}E)$ (Alice-Eve) as functions of QBER for both BB84 and B92. Identify the crossover point where $I(A{:}E) \geq I(A{:}B)$, beyond which no secret key can be extracted, and confirm this aligns with the GLLP-derived thresholds for each protocol.

    \item Determine and compare the maximum tolerable QBER thresholds for BB84 ($\sim$9.8\% with $f_{EC} = 1.16$) and B92 ($\sim$6.5\%), quantifying the practical noise tolerance gap between the two prepare and measure protocols.
\end{enumerate}

\subsection{Verification, Validation, and Benchmarking Visualisation}

\begin{enumerate}[label=\textbf{O5.\arabic*}, leftmargin=1.5cm]
    \item \textbf{Verification against theoretical predictions.} Confirm that simulation outputs match known analytical results:
    \begin{itemize}
        \item BB84 sifting efficiency converges to 50\% ($\pm$2\%).
        \item Full intercept-resend attack on BB84 produces QBER of 25\% ($\pm$2\%).
        \item Secure key rate reaches zero at the GLLP-predicted QBER thresholds for each protocol.
        \item B92 sifting efficiency converges to 25\%.
    \end{itemize}
    These checks serve as pass/fail criteria confirming that the simulation correctly implements the underlying quantum mechanics and post-processing mathematics.

    % potential edit on where i may want to validate results
    \item \textbf{Validation against published experimental data.} Compare simulated key rate versus distance curves against published experimental QKD demonstrations, including long-distance BB84 results \cite{Boaron2018} and metropolitan network deployments reported in recent survey literature \cite{Zhang2025}. The simulation is not expected to reproduce experimental results exactly (real systems include imperfections beyond the scope of this model), but the trends, orders of magnitude, and maximum operational distances should be broadly consistent. Discrepancies are documented and attributed to specific modelling assumptions.

    \item \textbf{Benchmarking visualisation suite.} Develop a set of figures to present simulation results and protocol comparisons, including:
    \begin{itemize}
        \item QBER versus eavesdropper interception rate for each protocol.
        \item Secure key rate versus fibre distance under realistic attenuation and $f_{EC} = 1.16$.
        \item Noise model comparison showing QBER response across depolarising, bit-flip, amplitude damping, and phase damping channels at varying strengths.
        \item Mutual information crossover plot ($I(A{:}B)$ vs.\ $I(A{:}E)$ vs.\ QBER).
    \end{itemize}
\end{enumerate}

\subsection{Real-World Feasibility and Broader Context}

\begin{enumerate}[label=\textbf{O6.\arabic*}, leftmargin=1.5cm]
    \item Evaluate the practical viability of BB84 and B92 over metropolitan fibre links (10-50 km), identifying the distance and noise regimes in which each protocol sustains a positive secure key rate with $f_{EC} = 1.16$. Determine, for a given noise floor, the maximum deployment distance at which each protocol remains operational.

    \item Contextualise the simulation results within the broader landscape of quantum-safe security, including the complementary roles of QKD and post-quantum cryptography (NIST FIPS 203/204/205), the "harvest now, decrypt later" threat model motivating near-term QKD deployment, and the economic and infrastructure barriers to adoption across sectors such as Research \& Development, Government, Banking \& Finance, Critical Infrastructure, Healthcare, and Emerging Industries.
\end{enumerate}

% ═══════════════════════════════════════════════════════════════════════════════    
\chapter{Literature Review}
% ═══════════════════════════════════════════════════════════════════════════════

% TODO: Review of QKD theory, protocol evolution, real-world deployments,
% and positioning of this work relative to existing research.

% repurpose the intro section into this paragraph

\section{The Quantum Threat to Classical Cryptography}
% new
The security of public-key cryptosystems such as RSA, Diffie-Hellman, and ECC depends on the assumed intractability of integer factorisation and the discrete logarithm problem. Shor's algorithm \cite{Shor1997} undermines both assumptions, solving each in polynomial time on a quantum computer. While no machine of sufficient scale exists today, resource estimates have fallen sharply. Gidney and Eker\r{a}~\cite{Gidney2021} showed that a 2048-bit RSA modulus could be factored in approximately eight hours using 20 million noisy qubits, and a 2025 follow-up by Gidney \cite{Gidney2025} reduced the requirement to under one million noisy qubits at the cost of roughly one week of runtime. As illustrated in Fig. \ref{fig:quantumAdvantage}, these declining estimates suggest we are approaching Quantum Economic Advantage (QEA): the point at which the practical resource cost of such an attack becomes feasible.

\begin{figure}[h]
    \centering
    \includegraphics[width=1\linewidth]{images.lit/quantumAdvantage.png}
    \caption{The Quantum Economic Advantage}
    \label{fig:quantumAdvantage}
\end{figure}

Symmetric-key cryptography is not immune. Grover's search algorithm \cite{Grover1996} provides a quadratic speed-up over brute-force key search, effectively halving the security level of schemes such as AES-128 and AES-256 \cite{QKDbible}. In practice, this means AES-128 offers only 64 bits of security against a quantum adversary, making AES-256 the minimum acceptable key length in a post-quantum setting. The asymmetry of impact is notable, where Shor's algorithm breaks public-key schemes outright, Grover's merely weakens symmetric schemes to a degree that can be compensated by doubling key sizes. It is for this reason that the quantum threat is primarily directed at public-key infrastructure.

The "harvest now, decrypt later" attack model gives these theoretical vulnerabilities immediate practical relevance. An adversary can intercept and store classically encrypted traffic today with the intent of decrypting it once a capable QC becomes available \cite{NIST}. For data with long confidentiality requirements, such as state secrets, medical records, and financial instruments, the effective window of vulnerability is already open. This threat model does not require a quantum computer to exist now; it simply requires that one will exist within the data's confidentiality lifetime.

Two complementary approaches address this threat. PQC replaces vulnerable mathematical assumptions with problems believed to resist quantum attack. NIST finalised its first three PQC standards in August 2024 (FIPS 203, 204, 205), with a fourth algorithm selected in March 2025 \cite{NIST_IR8547}. The accompanying IR-8547 transition roadmap deprecates quantum-vulnerable algorithms by 2030 and disallows them entirely by 2035.

However, as PQC provides \textit{computational} security, its guarantees only hold as long as the underlying mathematical problems remain hard, an assumption that cannot be proven and could be invalidated by future algorithmic breakthroughs. QKD takes a fundamentally different approach. By encoding key material in quantum states and exploiting physical principles such as the no-cloning theorem and measurement disturbance, QKD achieves \textit{information-theoretic} security, remaining secure regardless of an adversary's computational power, now or in the future \cite{Scarani2009, ShorPreskill2000}.

Shannon \cite{Shannon1945} established that information-theoretic secrecy requires a key at least as long as the message, a condition met by the one-time pad but impractical without a secure means of key distribution. QKD is the only known method of generating and distributing such keys at scale without relying on a trusted courier \cite{QKDbible}. Any eavesdropping attempt on the quantum channel inevitably disturbs the transmitted states, producing measurable errors that allow the communicating parties to detect interception and discard compromised key material. This physical guarantee is what distinguishes QKD from all classical and post-quantum alternatives. The role of QKD is therefore not as a competitor to PQC, but as a complementary layer offering a stronger class of security guarantee for applications where long-term confidentiality is critical.

% sources explaining the QM principles for how QKD works
\section{Quantum Mechanical Foundations}

To achieve information-theoretic security, every QKD protocol rests on a small set of quantum mechanical principles. This section introduces the five principles that are exploited by the BB84, B92, and E91 protocols. The treatment of these foundations follows Nielsen and Chuang \cite{QKDbible}, and the pedagogical review by \cite{Scarani2009}. 

\subsection{Qubits \& Superposition}
Consider a classical bit taking one of two definite values, 0 or 1. A Qubit (quantum bit) can exist in a linear combination of both 0 or 1, simultaneously. Using Dirac notation, the general single-qubit state is written as

%
\begin{equation}
\label{eq:qubit}
\ket{\psi} = \alpha\ket{0} + \beta\ket{1},
\end{equation}
%

where $\alpha$ and $\beta$ are complex probability amplitudes that satisfy $|\alpha|^{2}+|\beta|^{2}=1$ \cite{QKDbible}. The two states, $\ket{0}$ and $\ket{1}$ form the \textit{computational basis} (a.k.a. rectilinear or $Z$ basis). Another valid description of a qubit uses the diagonal basis (a.k.a $X$ Basis), defined as

%
\begin{equation}
\label{eq:diag-basis}
\ket{+} = \frac{\ket{0}+\ket{1}}{\sqrt{2}}, \qquad
\ket{-} = \frac{\ket{0}-\ket{1}}{\sqrt{2}}.
\end{equation}
%

In a physical QKD system, qubits are typically encoded in the polarisation of single photons: $\ket{0}$ and $\ket{1}$ may correspond to horizontal and vertical polarisation, while $\ket{+}$ and $\ket{-}$ correspond to $+45^{\circ}$ and $-45^{\circ}$ diagonal polarisation \cite{BB84}. Superposition is what allows a single photon to carry information in more than one basis at once, this property becomes a security asset because any attempt to read the qubit forces a basis choice that irreversibly disturbs the state.

\subsection{Measurement Disturbance}
Measuring a qubit that is in superposition collapses it probabilistically into one of the basis states of the chosen measurement apparatus. If Alice prepares $\ket{+}$ (encoding a logical 1 in BB84's diagonal basis) and Bob measures in the matching $X$ basis, he recovers the correct value with certainty. However, if Bob instead measures in the $Z$ basis, he obtains $\ket{0}$ or $\ket{1}$, each with probability $\frac{1}{2}$, and the original superposition is destroyed \cite{QKDbible}. The $Z$ and $X$ bases are called \textit{conjugate bases} because perfect knowledge in one precludes any information about the other. This is the physical mechanism behind the sifting step in BB84 and B92: only the rounds where Alice and Bob happen to use compatible bases contribute to the key, while mismatched rounds are discarded because Bob's result is uncorrelated with Alice's bit \cite{BB84, B92}.

Critically, measurement disturbance also applies to an eavesdropper (Eve). If Eve intercepts a qubit and measures it in the wrong basis, she collapses the state and re-sends a qubit that no longer matches Alice's preparation. When Alice and Bob later compare a random subset of their sifted key, they detect these errors as an elevated Quantum Bit Error Rate (QBER). An intercept-resend attack on BB84 introduces a characteristic QBER of approximately $25\%$ \cite{Scarani2009, ShorPreskill2000}. % NEW well above the $11\%$ threshold at which no secret key can be


\subsection{No-cloning Theorem}

A natural question is whether Eve could copy the qubit in transit, keep the copy for later analysis, and forward the original undisturbed qubit. The no-cloning theorem, proven independently by Wootters \& Zurek~\cite{WZ1982} and by Dieks \cite{Dieks1982}, rules this attack out. The theorem states that no physical process can produce a perfect copy of an arbitrary unknown quantum state. Formally, there exists no unitary operator $U$ such that $U\ket{\psi}\ket{0} = \ket{\psi}\ket{\psi}$, for all $\ket{\psi}$ \cite{QKDbible}. The proof is compact: suppose such a $U$ did exist for two non-orthogonal states $\ket{\psi}$ and $\ket{\phi}$. Unitarity requires preservation of inner products, giving $\braket{\psi|\phi} = (\braket{\psi|\phi})^{2}$, which is satisfied only when $\braket{\psi|\phi} = 0$ or $1$. Since $\ket{\psi}$ and $\ket{\phi}$ are non-orthogonal by assumption, neither condition holds, producing a contradiction.

No-cloning is the single most important result underpinning QKD. It guarantees that Eve cannot silently duplicate quantum transmissions and is the reason quantum key distribution can offer information-theoretic security rather than the computational security of classical schemes~\cite{ekertQI}. B92 exploits this principle particularly directly, where the protocol encodes bits using only two non-orthogonal states ($\ket{0}$ and $\ket{+}$), no measurement or cloning strategy can distinguish them with certainty \cite{B92}.

\subsection{Heisenberg's Uncertainty}

Closely related to measurement disturbance is Heisenberg's uncertainty principle \cite{Heisenberg1927}. In its information-theoretic form relevant to QKD, the principle states that for two non-commuting observables (such as polarisation measured in
the $Z$ and $X$ bases), simultaneously obtaining precise values for both is impossible. Any measurement that yields information about one observable necessarily introduces uncertainty in the other \cite{QKDbible}. In the context of QKD, this means Eve faces a fundamental trade-off: gaining information about Alice's bit in one basis unavoidably disturbs the conjugate-basis encoding, which manifests as detectable errors. Heisenberg's principle and the no-cloning theorem are logically distinct but operationally complementary; together they establish that any eavesdropping strategy, whether based on measurement, cloning, or a combination of both, will leave a detectable footprint in the quantum channel \cite{Scarani2009}. 

\subsection{Quantum Entanglement \& Bell's Theorem}

Two or more qubits can exist in an \textit{entangled} state that cannot be described as a product of individual qubit states. The canonical example is the Bell state (also called an EPR pair after Einstein, Podolsky and Rosen \cite{EPR1935}): % NEW
%
\begin{equation}
\label{eq:bell-state}
\ket{\Phi^{+}} = \frac{\ket{00}+\ket{11}}{\sqrt{2}}.
\end{equation}
%
If Alice holds the first qubit and Bob the second, a measurement by Alice in the $Z$ basis yielding $\ket{0}$ instantaneously projects Bob's qubit into $\ket{0}$ as well, regardless of the spatial separation between them~\cite{QKDbible}. Entanglement does not permit faster-than-light communication (the individual outcomes are still uniformly random), but it does produce correlations that cannot be explained by any local hidden-variable theory.

Bell \cite{Bell1964} formalised this distinction in 1964 by deriving an inequality that any local realist theory must satisfy. The experimentally testable form due to Clauser, Horne, Shimony and Holt \cite{CHSH1969} (the CHSH inequality) states
%
\begin{equation}
\label{eq:chsh}
|S| = |E(a,b) - E(a,b') + E(a',b) + E(a',b')| \leq 2,
\end{equation}
%
where $E(a,b)$ denotes the correlation of measurement outcomes for settings $a$ and $b$. Quantum mechanics predicts a maximum value of $|S| = 2\sqrt{2} \approx 2.83$, a violation confirmed experimentally by Aspect \cite{Aspect1982}.

Ekert's E91 protocol \cite{E91} directly leverages Bell inequality violation as a built-in eavesdropping test. While BB84 and B92 do not require entanglement, the Shor-Preskill security proof \cite{ShorPreskill2000} establishes a formal equivalence between BB84 and an entanglement-based scheme, meaning the security of prepare-and-measure protocols ultimately derives from the same non-local correlations.

\newpage
% QKD whitepapers
\section{Quantum Key Distribution Protocols}
\begin{figure}[h]
    \centering
    \includegraphics[width=0.6
\linewidth]{images.lit/QKDclassification.png}
    \caption{Classification of QKD Protocol Schemes}
    \label{fig:QKDclassification}
\end{figure}


QKD Protocols can be classified into four basic aspects of use, as depicted in amber in Fig. \ref{fig:QKDclassification}. Each aspect is then further classified based upon the quantum mechanical principles exploited, seen in green. Continuous  Variable (CV) schemes encode data using the continuous variable of a quantum system, such as the amplitude or phase of a light wave. Due to this, CV protocols can be used with common optical equipment and, as such, are less susceptible to noise when deployed practically. In comparison, Discrete Variable (DV) schemes are implemented using discrete quantum states, such as single polarized photons or entangled-pairs of photons \cite{Kumar2025}. Furthermore, QKD protocol structures can classified into three broad paradigms: Prepare \& Measure, Entanglement-Based, or Measurement-Device-Independent schemes. This project focuses on three DV-Protocols: BB84, B92 P\&M protocols, and the Entanglement-Based E91 protocol. 

% edit below this
\subsection{BB84}
\label{sec:bb84}
\begin{figure}[ht]
    \centering
    \includegraphics[width=0.6\linewidth]{images.lit/QuantumAbstract.png}
    \caption{Abstract P\&M QKD Network Model}
    \label{fig:QKDnetworkmodel}
\end{figure}

The BB84 protocol, proposed by Bennett and Brassard in 1984 \cite{BB84}, was the first practical QKD scheme. It is a prepare-and-measure protocol that encodes classical bit values onto single-photon polarisation states drawn from two mutually unbiased bases: the rectilinear basis ($\ket{0},\ket{1}$) and the diagonal basis ($\ket{+},\ket{-}$). The sender, Alice, randomly selects both a bit value and a basis for each qubit, producing one of four possible states:
%
\begin{equation}
  \ket{0},\quad \ket{1},\quad
  \ket{+} = \frac{\ket{0}+\ket{1}}{\sqrt{2}},\quad
  \ket{-} = \frac{\ket{0}-\ket{1}}{\sqrt{2}}.
  \label{eq:bb84-states}
\end{equation}
%


These states are transmitted to Bob over the quantum channel. Bob then independently and randomly chooses a measurement basis for each incoming qubit. Because the two bases are mutually unbiased, a measurement in the wrong basis yields a uniformly random outcome and provides no information about Alice's bit value \cite{QKDbible}.

After transmission is complete, Alice and Bob switch to an authenticated classical channel to perform \textit{sifting}, as seen in Fig. \ref{fig:QKDnetworkmodel}. At this point, Bob announces only his choice of basis for each qubit, and Alice confirms which choices matched her preparation basis. All rounds in which the bases disagree are discarded. In an ideal, noiseless channel, the remaining \textit{sifted key} is perfectly correlated between Alice and Bob, and roughly half the transmitted qubits survive the sifting step \cite{BB84}.

\begin{figure}[h]
    \centering
    \includegraphics[width=0.7\linewidth]{images.lit/bb84operation.png}
    \caption{BB84 Operation Model}
    \label{fig:BB84operation}
\end{figure}

To check for eavesdropping, Alice and Bob sacrifice a random subset of the sifted key and compare bit values over the public channel. Any intercept-resend attack by an eavesdropper, Eve, necessarily disturbs a fraction of the qubits because she cannot simultaneously determine the polarisation in both bases without violating the no-cloning theorem \cite{QKDbible}. A full intercept-resend strategy is expected to produce a Quantum Bit Error Rate (QBER) of approximately 25\%, well above the theoretical security threshold of approximately 11\% \cite{ekertQI}. If the observed QBER exceeds this threshold, the protocol aborts. Otherwise, Alice and Bob proceed through error correction and privacy amplification to yield a final secret key, depicted in Fig. \ref{fig:BB84operation}.


\subsection{B92}
\label{sec:b92}

\begin{figure}[h]
    \centering
    \includegraphics[width=0.6\linewidth]{images.lit/b92basis.png}
    \caption{B92 Measurements Bases}
    \label{fig:B92bases}
\end{figure}

In 1992, Bennett proposed the B92 protocol as a simplification of BB84 \cite{B92}. The central insight is that unconditional security does not require four encoding states; two non-orthogonal states (seen Fig. \ref{fig:B92bases} are sufficient because no measurement can perfectly distinguish between non-orthogonal quantum states \cite{ekertQI}. In B92, Alice encodes bit value 0 as $\ket{0}$ and bit value 1 as $\ket{+}$:

%
\begin{equation}
  \ket{\varphi} =
  \begin{cases}
    \ket{0}  & \text{if } a = 0, \\[4pt]
    \ket{+}  & \text{if } a = 1.
  \end{cases}
  \label{eq:b92-states}
\end{equation}
%

Bob generates his own random bit $a'$ and measures in the $Z$ basis ($\ket{0},\ket{1}$) if $a' = 0$, or in the $X$ basis ($\ket{+},\ket{-}$) if $a' = 1$. When Bob's measurement basis is incompatible with Alice's state, the outcome is conclusive only some of the time, hence, inconclusive results (a.k.a. erasures) are discarded \cite{QKDbible}. Bob publicly announces which rounds yielded a conclusive result, and Alice retains only the corresponding bits. Neither party reveals their basis choice or bit value during this exchange.

Because B92 relies on only two states rather than four, its sifting efficiency is lower than that of BB84, and its tolerance to channel noise is reduced. The theoretical QBER security threshold for B92 is approximately 6.5\%, compared to approximately 11\% for BB84 \cite{ekertQI}. Furthermore, the two-state encoding makes B92 more susceptible to unambiguous state discrimination (USD) attacks, in which Eve exploits the known non-orthogonal structure to occasionally identify Alice's state with certainty while flagging ambiguous cases as losses \cite{B92}. Despite these limitations, B92 remains an important protocol both for its theoretical elegance and because it demonstrates that the minimal resource needed for secure QKD is a pair of non-orthogonal states.

\subsection{E91}
\label{sec:e91}

In contrast to the P\&M paradigm of BB84 and B92, the E91 protocol proposed by Ekert in 1991 \cite{E91} derives its security from quantum entanglement and Bell's theorem. A source prepares pairs of qubits in the singlet Bell state:

%
\begin{equation}
  \ket{\Psi^{-}} = \frac{\ket{01} - \ket{10}}{\sqrt{2}},
  \label{eq:e91-singlet}
\end{equation}
%

and distributes one qubit from each pair to Alice and the other to Bob. Each party independently selects a random measurement orientation from a predetermined set of angles. When Alice and Bob happen to choose the same measurement basis, their outcomes are perfectly anti-correlated, and these matching rounds are used to construct the raw key \cite{E91}.

The rounds in which Alice and Bob chose different bases serve a separate purpose: they are used to evaluate the Clauser-Horne-Shimony-Holt (CHSH) inequality \cite{QKDbible}. For the singlet state, quantum mechanics predicts a CHSH parameter of $S = 2\sqrt{2}$, which maximally violates the classical bound of $S \leq 2$. Any eavesdropping that disturbs the entanglement between the pairs will reduce the observed value of $S$ towards or below the classical limit, thereby revealing Eve's presence \cite{ekertQI, E91}. This connection between eavesdropper detection and Bell inequality violation is the defining feature of E91 and distinguishes it fundamentally from P\&M protocols.

This entanglement-based approach offers a conceptual advantage as the key bits are not determined until the moment of measurement, so there is no point during transmission at which an encoded classical value exists for Eve to intercept. However, E91 places greater demands on hardware because it requires a source of high-fidelity entangled pairs and the ability to distribute them over the quantum channel while preserving entanglement \cite{ekertQI}. In a simulation context, this translates to the additional complexity of modelling entangled state preparation, Bell-state measurement, and CHSH parameter estimation alongside the standard post-processing pipeline shared with BB84 and B92.



% discuss the basis sifitng & priavcy amplification processes
\section{The Post-Processing Pipeline}
Once the quantum transmission phase of a QKD protocol is complete, the raw measurement outcomes held by Alice and Bob must be transformed into a shared secret key through a sequence of classical post-processing steps. This pipeline is common to all three protocols considered in this project, although the details of sifting differ between prepare-and-measure and entanglement-based schemes.

The first stage is \textit{sifting}, or basis reconciliation. In BB84, Alice and Bob publicly compare their basis choices and discard all rounds where their bases did not match, retaining approximately 50\% of transmitted qubits \cite{BB84}. However, in B92, Bob announces which rounds yielded a conclusive measurement outcome, and only those rounds are kept, producing a sifting efficiency of approximately 25\% \cite{B92}. For E91, sifting separates the rounds into two groups based on whether Alice and Bob selected the same measurement angle. Matching-angle rounds contribute to the raw key, while non-matching rounds are reserved for CHSH parameter estimation \cite{E91}.

Following sifting, Alice and Bob perform \textit{error estimation} by publicly comparing a randomly chosen subset of the sifted key. The fraction of disagreements in this subset provides an estimate of the QBER. If the observed QBER exceeds the protocol-specific security threshold, the key is deemed compromised and the protocol aborts. Otherwise, the sacrificed bits are discarded and the remaining sifted key proceeds to error correction \cite{Scarani2009}.

\textit{Error correction} reconciles the remaining discrepancies between Alice's and Bob's key strings. The CASCADE protocol, introduced by Brassard and Salvail \cite{Cascade1994}, is the most widely used scheme in QKD implementations. CASCADE operates through multiple passes of binary search, where in each pass, the key is divided into blocks of increasing size, and parity checks identify and correct single-bit errors within each block. The multi-pass structure progressively eliminates errors that were missed in earlier rounds. A critical parameter is the error correction efficiency $f_{EC}$, defined as the ratio of information actually leaked during correction to the Shannon-limit minimum. CASCADE achieves $f_{EC} \approx 1.16$ in practice \cite{Cascade1994}, meaning it leaks approximately 16\% more information than the theoretical minimum. This distinction has direct consequences for the achievable secure key rate, as discussed below.

The final stage is \textit{privacy amplification}, which compresses the corrected key to remove any information that may have leaked to an eavesdropper during both the quantum transmission and the error correction phases. Universal hashing functions are applied to map the corrected key of length $n$ to a shorter final key of length $\ell < n$, where $\ell$ is chosen such that Eve's mutual information with the final key is negligibly small \cite{Bennett1995PA}. The amount of compression required depends on both the estimated QBER and the information leaked during error correction.

The GLLP formula \cite{GLLP2004} quantifies how much key material survives this entire pipeline. For BB84 with a single-photon source, the asymptotic secure key rate per sifted bit is given by
%
\begin{equation}
\label{eq:gllp}
R = 1 - h(\text{QBER}) - f_{EC} \cdot h(\text{QBER}),
\end{equation}
%
where $h(\cdot)$ denotes the binary Shannon entropy. Setting $f_{EC} = 1.0$ (the Shannon limit) yields the well-known 11\% QBER threshold for BB84. However, using the realistic CASCADE efficiency of $f_{EC} = 1.16$ shifts this threshold to approximately 9.8\%, a reduction with meaningful engineering consequences for budget calculations. It is worth noting that many simulation studies in the literature default to $f_{EC} = 1.0$, which overestimates the achievable secure key rate. This project uses $f_{EC} = 1.16$ throughout to provide a more accurate representation of practical system performance.

\section{Noise Models for Modelling} 

In any practical QKD deployment, the quantum channel introduces errors that are physically indistinguishable from those caused by eavesdropping. Accurately modelling these channel imperfections is therefore essential for a simulation platform that seeks to evaluate protocol performance under realistic conditions.

The \textit{depolarising channel} is the most commonly used noise model in QKD analysis. It maps an arbitrary single-qubit state $\rho$ to the maximally mixed state with probability $p$, effectively applying a random Pauli error:
%
\begin{equation}
\label{eq:depolarising}
\mathcal{E}_{dep}(\rho) = (1-p)\rho + \frac{p}{3}(X\rho X + Y\rho Y + Z\rho Z).
\end{equation}
%
Because the depolarising channel applies errors isotropically across all axes of the Bloch sphere, it affects all encoding bases equally and is therefore compatible with the assumptions underlying the GLLP key rate formula \cite{GLLP2004}. Physically, it represents a worst-case model of isotropic noise arising from environmental decoherence in the fibre channel.

% block sphere image

The \textit{phase damping} channel models the loss of phase coherence without energy exchange, representing the physical process of photon scattering in optical fibre. Its action on a density matrix is:
%
\begin{equation}
\label{eq:phase-damping}
\mathcal{E}_{pd}(\rho) = (1-\lambda)\rho + \lambda Z\rho Z,
\end{equation}
%
where $\lambda$ is the damping probability. Phase damping introduces a basis-dependent asymmetry, acting as the identity on $Z$-basis eigenstates ($\ket{0}$, $\ket{1}$) while decohering superposition states in the $X$ basis ($\ket{+}$, $\ket{-}$). This asymmetry has protocol-differentiating consequences. BB84, which encodes in both bases symmetrically, experiences errors only on the $X$-basis fraction of its transmissions. B92, however, encodes bit value 1 as $\ket{+}$, a superposition state that is directly vulnerable to phase damping, while bit value 0 is encoded as $\ket{0}$, a $Z$-eigenstate that is unaffected \cite{QKDbible}. This creates a bit-value-dependent error rate unique to B92 under phase noise.

The \textit{amplitude damping} channel models energy dissipation, physically corresponding to photon absorption in the fibre medium. It is characterised by a damping parameter $\gamma$ that governs the probability of a $\ket{1} \rightarrow \ket{0}$ transition. In the context of QKD, amplitude damping provides an approximation of fibre loss at the qubit level, although true photon loss is more accurately described as a transmission efficiency $\eta$ applied at the classical level. This distinction is discussed further in the methodology.

Several other noise channels available in Qiskit Aer were considered and excluded. The bit-flip channel is a special case subsumed by depolarising noise. Thermal relaxation requires $T_1$/$T_2$ time constants that are not physically meaningful for single-photon transmission through fibre. The reset channel models spontaneous emission to the ground state, a process not relevant to photonic QKD systems. This selection prioritises physical relevance to fibre-optic QKD over exhaustive coverage of the models available in Qiskit Aer.


\section{QKD Simulation Frameworks}

Several software frameworks exist for simulating quantum communication protocols, each with distinct design philosophies and trade-offs relevant to the goals of this project.

Qiskit Aer \cite{QiskitAer}, the noise simulation backend of IBM's Qiskit framework, provides gate-level circuit simulation with a comprehensive built-in noise model library. Its \texttt{NoiseModel} API allows noise channels to be attached to specific gates rather than applied globally, enabling the channel-marker approach used in this project: an identity gate is inserted between state preparation and measurement to represent the quantum channel, and noise is applied exclusively to this gate. This preserves the trusted-device assumption by ensuring that Alice's preparation and Bob's measurement remain noise-free. No other framework surveyed offered this level of noise insertion control with comparable accessibility. Qiskit Aer's Python-native interface and active maintenance by IBM further supported its selection as the simulation backend for this project.

% include a comment on Pennylane
Pennylane...

\section{Research Gaps}
My aim to present the protocols under unified modelling constraints, to obtain a critical analysis of BB84, B92, and E91.

% fix references
The existing QKD simulation literature exhibits several gaps that this project aims to address. The majority of published simulation studies focus on benchmarking a single protocol in isolation. For example, Ma \textit{et al.} \cite{Ma2005} analyse BB84 key rates under various source assumptions, while Tamaki \textit{et al.} \cite{Tamaki2014} examine B92 security bounds independently. Comparative studies that evaluate multiple protocols under identical channel and noise conditions on the same simulation platform are considerably less common, and those that do exist typically compare only two protocols or restrict themselves to a single noise model.

Furthermore, few existing simulation platforms compare prepare-and-measure and entanglement-based protocols side-by-side using a unified benchmarking framework. BB84 and B92 are frequently analysed together due to their shared P\&M structure, but direct comparisons with E91 under matched noise conditions are rare. This limits the ability to draw meaningful conclusions about the relative strengths and weaknesses of these two fundamentally different QKD paradigms.

A further gap concerns the treatment of post-processing parameters in simulation studies. Many published analyses assume Shannon-limit error correction ($f_{EC} = 1.0$), which produces optimistic key rate estimates and overstates the maximum tolerable QBER. In practice, the CASCADE protocol operates at $f_{EC} \approx 1.16$ \cite{Cascade1994}, and this difference has direct consequences for the engineering viability of a given fibre link. Simulations that neglect this factor risk misinforming deployment decisions.

This project addresses these gaps by providing a simulation platform that benchmarks BB84, B92, and E91 under identical, configurable noise conditions with GLLP-based key rate analysis using realistic error correction efficiency. The platform is designed to enable fair comparison between P\&M and entanglement-based QKD protocols on a unified framework.




% ═══════════════════════════════════════════════════════════════════════════════
\chapter{Methodology}
% ═══════════════════════════════════════════════════════════════════════════════

\section{Simulation Architecture}

% TODO: Describe the modular object-oriented design of the platform.

\section{Protocol Implementations}

% --- HEADING LEVEL DEMO ---
% These demonstrate all four IEEE heading levels within a section:

\subsection{BB84 Protocol}

% Heading 3 example (run-in, indented 1em, Arabic numeral + right paren):
\subsubsection{Basis Selection} The sender (Alice) randomly selects
between the rectilinear and diagonal bases for each qubit transmission.

% Heading 4 example (run-in, indented 2em, lowercase alpha + right paren):
\paragraph{Rectilinear basis} Uses the computational basis states
$\ket{0}$ and $\ket{1}$ for encoding bit values 0 and 1 respectively.

\paragraph{Diagonal basis} Uses the Hadamard-rotated states
$\ket{+}$ and $\ket{-}$ for encoding.

% --- Example: Code Listing ---
\Cref{lst:measure} shows the measurement logic used by Bob
in the BB84 protocol implementation.

\begin{lstlisting}[language=Python, caption={BB84 Bob's measurement logic}, label={lst:measure}]
def measure_qubits(qubits, bases):
    """
    Measurement logic for QKD receiver.
    Applies rectilinear or diagonal measurement
    depending on the randomly chosen basis.
    """
    results = []
    for q, b in zip(qubits, bases):
        if b == 0:
            results.append(measure_rectilinear(q))
        else:
            results.append(measure_diagonal(q))
    return results
\end{lstlisting}

% --- Example: Algorithm ---
The high-level flow of the BB84 key sifting procedure is presented in
\cref{alg:bb84}.

\begin{algorithm}[H]
    \caption{BB84 Key Sifting Procedure}\label{alg:bb84}
    \KwIn{Alice's bits $a_i$, bases $\hat{a}_i$; Bob's measurements $b_i$, bases $\hat{b}_i$}
    \KwOut{Sifted key $K$}
    $K \leftarrow \emptyset$\;
    \For{$i \leftarrow 1$ \KwTo $n$}{
        Alice publishes $\hat{a}_i$ over classical channel\;
        Bob publishes $\hat{b}_i$ over classical channel\;
        \If{$\hat{a}_i = \hat{b}_i$}{
            Append $b_i$ to $K$ \Comment{Matching basis - keep bit}
        } 
    }
    \Return{$K$}
\end{algorithm}

\subsection{B92 Protocol}

% TODO: Implementation details for B92.

\subsection{E91 Protocol}

% TODO: Implementation details for E91 (entanglement-based).

%modelling assumptions for the protocols: measurement device characteristics etc
\section{Modelling Assumptions}

\section{Noise Models and Channel Simulation}

% TODO: Depolarising channel, amplitude damping, eavesdropper models.

\section{Performance Metrics}

% TODO: Secure key rate (Shor--Preskill bound, GLLP), QBER, sifted key
% efficiency, error correction overhead ($f_{\mathrm{EC}} = 1.16$).

% --- Example: Equation ---
% IEEE equations: 10pt, 12pt line spacing, centred, optional right-justified numbering
The asymptotic secure key rate under the Shor--Preskill bound is given by:
\begin{equation}\label{eq:keyrate}
    \keyrate = 1 - h(\QBER) - f_{\mathrm{EC}}\, h(\QBER)
\end{equation}
where $h(\cdot)$ denotes the binary Shannon entropy and $f_{\mathrm{EC}}$
is the error correction inefficiency factor.

% --- Example: Table ---
% IEEE tables: caption above, 8pt caption, Roman numeral, small caps "TABLE"



\chapter{Results \& Discussion}

% TODO: Present and analyse simulation results.
% Include comparative analysis of BB84, B92, E91 under varying conditions.

% simulation outputs
\section{Simulation Outputs}

% bb84 vs b92
\section{BB84 vs B92}

% performance validation
\section{Performance Validation}

% implications of results
\section{Implications}

% ═══════════════════════════════════════════════════════════════════════════════
\chapter{Conclusion \& Future Work}
% ═══════════════════════════════════════════════════════════════════════════════

% TODO: Summarise findings, contributions, limitations, and future directions.




% ─────────────────────────────────────────────────────────────────────────────
% REFERENCES
% ─────────────────────────────────────────────────────────────────────────────
% IEEE references: 8pt, 9pt line spacing, bracketed Arabic numerals [1], [2]
\clearpage
{\footnotesize
\bibliographystyle{IEEEtran}
\bibliography{bib}
}

% ─────────────────────────────────────────────────────────────────────────────
% APPENDICES
% ─────────────────────────────────────────────────────────────────────────────
\appendix

\chapter{Source Code}
% TODO: Full source code listings or link to repository.

\chapter{Additional Results}
% TODO: Supplementary plots, extended data tables.

\end{document}